\chapter{Arguments, exponentielle imaginaire et trigonométrie}
\section*{Exercice 26.1 -- Forme trigonométrique}
Le module vaut $2$ et un argument est $2\pi/3$, car le cosinus vaut $-1/2$ et le sinus $\sqrt3/2$. Ainsi $z=2(\cos(2\pi/3)+i\sin(2\pi/3))=2e^{2i\pi/3}$. Les arguments sont $2\pi/3+2k\pi$.

\section*{Exercice 26.2 -- Puissance}
$\sqrt3+i=2e^{i\pi/6}$. La formule de Moivre donne la puissance $2^6e^{i\pi}=-64$. Les modules se multiplient et les arguments s'additionnent.

\section*{Exercice 26.3 -- Addition trigonométrique}
Développer $(\cos a+i\sin a)(\cos b+i\sin b)$ puis identifier les parties réelle et imaginaire donne $\cos(a+b)=\cos a\cos b-\sin a\sin b$ et $\sin(a+b)=\sin a\cos b+\cos a\sin b$. Dans le cours, le produit exponentiel est lui-même établi géométriquement avant cette réécriture.

\section*{Exercice 26.4 -- Linéarisation}
Avec $2\cos x=e^{i x}+e^{-i x}$, le cube donne $8\cos^3x=e^{3i x}+e^{-3i x}+3(e^{i x}+e^{-i x})=2\cos3x+6\cos x$. Donc $\cos^3x=(\cos3x+3\cos x)/4$.

\section*{Exercice 26.5 -- Somme de phases}
C'est une somme géométrique : $(1-e^{i nt})/(1-e^{i t})$. Si $t$ est multiple de $2\pi$, tous les termes valent $1$ et la somme vaut $n$ ; la formule quotient n'est pas définie et ne doit pas être utilisée directement.

\section*{Exercice 26.6 -- Argument d’un produit}
La somme $3\pi/2$ est un argument, mais elle n'appartient pas à l'intervalle choisi. Soustraire $2\pi$ donne l'argument principal $-\pi/2$. Les égalités d'arguments sont des congruences modulo $2\pi$.
